% =====================================================================
%  第 4 章  如何判断做得好不好（草稿）
% =====================================================================

\section{如何判断做得好不好}
\label{sec:evaluation}

\begin{notebox}[本章状态：草稿]
本章描述评价方案。相关实验尚未进行，因此没有结果。下文的指标定义会在实现后根据实际情况调整。
\end{notebox}

\subsection{为什么“发现了新猜想”不能作为主要指标}

最直觉的指标当然是“它提出了多少条后来被证明的猜想”。但这一条在课程周期内根本无法评估：一条猜想从提出到被证明往往需要数年，我们等不起。

退一步，可以数“提出了多少条新猜想”。这同样不可靠：新不新需要查文献，而“看起来没人提过”往往只是我们没查到。更糟的是，这个指标会诱导系统去追求古怪而无法验证的形式——越冷门越“新”。

因此我们采用一个更保守但可操作的标准。

\subsection{主要判据：一小时可判真伪}

我们采用的主要判据是：\textbf{系统提出的猜想，能否让一个懂图论的人在不到一小时内判断出真假}。

这个标准的好处是它同时惩罚了两类失败：如果提出的猜想太弱（比如只是已知定理的改写），人一眼就能看出，但那不构成发现；如果提出的猜想太怪（比如涉及没人研究过的量），人根本无从下手，这也不算有用的产出。

按这个标准，一条“好”的猜想应该是：形式简单、涉及的量都是常见的、在数据上高度可信、且现有的反例搜索没能推翻它。

\subsection{辅助判据：重新发现已知结论}

在系统正式运行之前，我们先让它去“发现”那些我们已经知道答案的定理——比如式 \eqref{eq:chi-alpha}。

这一项是\textbf{正确性检验}，不是成果展示。如果系统连这些显然成立的不等式都找不到，或者找出来又自己推翻了，那说明流程里有 bug，后面所有的结果都不可信。反过来说，系统能重新发现这些结论，只能说明它没坏，不能说明它有用。

这个区分很重要：汇报时如果把“重新发现了 $\chi+\alpha\le n+1$”当成主要成果来讲，会显得很空洞。它的正确定位是“系统自检通过”。

\subsection{对照实验}

为了说明“统计筛选”这个设计是否真的有用，我们计划做一组对照：同一批候选，一半经过多重检验校正，一半不经过，比较两者最终保留集合里假货的比例。

判断“假货”需要一个标准答案，我们用两种方式构造：一是人为注入已知为假的不等式（比如把某个定理的不等号反向），看系统能否把它们筛掉；二是在更大的图上重新验证，把在小图上成立、在大图上失效的候选标记出来。

如果校正后的假货比例明显更低，说明设计有效；如果没有差别，说明这个环节可以去掉——这同样是有价值的结论，能省下后来者的功夫。

\subsection{要报告的负面结果}

课程作业容易只报喜不报忧，我们打算反过来，明确列出几类会如实报告的负面结果：

\begin{itemize}
  \item \textbf{LLM 的贡献有限}。如果消融实验显示 LLM 与随机选择无显著差异，就照实写。
  \item \textbf{统计筛选的代价过高}。多重检验校正会降低召回率，如果代价大到“筛完之后什么都不剩”，那这个方法在实践上就没用。
  \item \textbf{系统只能在小图上工作}。如果所有结论都局限在七八个顶点，那它的数学意义就很有限，这一点必须写明。
\end{itemize}

这些结果不会让课题显得失败。相反，一个能清楚说明“这个方法在什么范围内有效、在什么范围内失效”的报告，比一个只展示成功案例的报告更有说服力。
